\chapter{Research Novelty}
\label{chap:novelty}

Before diving into the numerical machinery, it is worth stating
plainly \emph{what is new} about this study.  The lunar thermal
modelling literature is mature --- \citet{wesselink1948,cremers1971,
vasavada1999,mitchell2008,hayne2017,schorghofer2020} have each
contributed a cornerstone result.  Yet none of these works combine, in
a single open pipeline, the four ingredients listed below.  Each
ingredient is load-bearing for the science target; removing any one of
them materially weakens the conclusions.

\section{Multi-site blind validation against every in-situ record}
\label{sec:novelty-multisite}

Most published lunar thermal solvers are validated against a
\emph{single} data set --- typically either Apollo~15 or Diviner ---
and are then extrapolated to new geometries without a second
independent check.  The TSUKIMI pipeline, by contrast, is regression-
tested against the complete set of in-situ subsurface thermal records
currently available in the peer-reviewed literature
(\cref{tab:validation-sites}).

\begin{table}[H]
\centering
\small
\begin{tabularx}{\linewidth}{l l l Y}
\toprule
Site & Latitude & Instrument & Reference \\
\midrule
Apollo~15  & \SI{26}{\degree}N & HFE thermocouple probes       & \citet{langseth1976,nagihara2018} \\
Apollo~17  & \SI{20}{\degree}N & HFE thermocouple probes       & \citet{langseth1976,nagihara2018} \\
Chang'E-4  & \SI{45}{\degree}S & YuTu-2 surface thermistors    & \citet{huang2022} \\
ChaSTE     & \SI{69}{\degree}S & Chandrayaan-3 penetrometer    & \citet{murty2025,das2025,seth2025} \\
\bottomrule
\end{tabularx}
\caption[Validation sites]{The four in-situ lunar thermal data sets
  currently available; the TSUKIMI pipeline is validated against all
  four simultaneously using the same solver kernel.}
\label{tab:validation-sites}
\end{table}

\begin{novelty}
No single study to date has \emph{blindly} reproduced all four of
these records with one unchanged solver.  The only parameter that
varies between sites is the H-scale $\Hscale$ (or equivalently the
discrete layer thickness) --- the physics kernel is identical.  This is
the ``cross-validation'' protocol that would be demanded of any
terrestrial climate model and it removes the usual degree of freedom
where an author can tune a bespoke parameter set to each landing site.
\end{novelty}

\section{Ice-coupled conductivity feedback}
\label{sec:novelty-ice}

In the polar night and in permanently shadowed regions, water ice
becomes thermodynamically stable within the upper ${\sim}\SI{1}{\meter}$
of regolith \citep{paige2010,schorghofer2020}.  The thermal
conductivity of hexagonal water ice follows Klinger's relation
\citep{klinger1980}
\begin{equation}
  \Kice(T) = \frac{567}{T}\ \si{\watt\per\meter\per\kelvin},
  \qquad T \in [\SI{20}{\kelvin},\SI{273}{\kelvin}],
  \label{eq:klinger}
\end{equation}
which at $T=\SI{100}{\kelvin}$ yields
$\Kice = \SI{5.67}{\watt\per\meter\per\kelvin}$ --- \emph{three orders
of magnitude} larger than the dry-regolith value
($\Kreg \sim \SI{1e-3}{\watt\per\meter\per\kelvin}$).  The effective
conductivity of an ice--regolith mixture at volumetric ice fraction
$\phi$ is dominated by the ice phase for any $\phi \gtrsim 0.1$:
\begin{equation}
  k_{\mathrm{eff}}(T,\phi) =
  \phi\,\Kice(T) + (1-\phi)\,\Kreg(T,z),
  \label{eq:k_mix}
\end{equation}
which is a non-trivial function of \emph{both} temperature and depth.

\begin{concept}
Ice-coupled conductivity creates a positive feedback: once ice forms, it
raises $k_{\mathrm{eff}}$, which raises the diurnal damping depth,
which stabilises ice against the next sublimation cycle.  A solver that
ignores this loop (i.e.\ fixes $k$ at the dry-regolith value) will
systematically under-predict the ice-stability depth by tens of
centimetres.
\end{concept}

\begin{novelty}
The TSUKIMI pipeline is --- to the author's knowledge --- the first
open-source lunar thermal code to close this loop self-consistently
within the time-stepping kernel, rather than applying it as a
post-processing correction.  The implementation is the subject of
\cref{chap:phase3}.
\end{novelty}

\section{Pixel-level topographic slope correction}
\label{sec:novelty-slope}

ChaSTE landed at \SI{69}{\degree}S on the rim of a minor crater with a
local slope estimated at \SI{6}{\degree} from the Chandrayaan-3
descent imagery \citep{seth2025}.  A flat-surface thermal model
under-predicts the observed peak surface temperature at this site by
approximately \SI{40}{\kelvin} (\cref{chap:phase1}, §\,ChaSTE).  This
bias is not a physics failure of the solver --- it is a geometry
failure of the boundary condition: a flat surface receives
systematically less noon-time insolation at \SI{69}{\degree}S than the
pole-facing slope on which the lander actually sits.

The corrected cosine-of-incidence is
\begin{equation}
  \cos i \;=\; \sin\beta\,\sin s\,\cos(\alpha - \varphi_s)
            \;+\; \cos\beta\,\cos s,
  \label{eq:cos_incidence}
\end{equation}
where $\beta$ is the solar zenith angle, $s$ the local slope, $\alpha$
the solar azimuth, and $\varphi_s$ the slope azimuth (aspect).
Applying \cref{eq:cos_incidence} with $s=\SI{6}{\degree}$ and a
pole-facing aspect recovers the observed peak temperature within
\SI{10}{\kelvin}.

\begin{novelty}
Most published polar models apply slope correction only in aggregate
(\emph{climate-zone} averages).  TSUKIMI applies \cref{eq:cos_incidence}
\emph{per pixel}, using the LOLA/Kaguya merged DEM as the slope source,
which is the required granularity to compare against point
measurements like ChaSTE.
\end{novelty}

\section{End-to-end reproducible workflow}
\label{sec:novelty-reproducibility}

The last novelty is less scientific than methodological but equally
important for a PhD-scale study.  Every notebook in the repository is
\emph{self-installing} and \emph{self-fetching}:

\begin{lstlisting}[caption={Bootstrap cell present at the top of every notebook.},label={lst:bootstrap}]
# --- Auto-install + auto-fetch -------------------------------------
import subprocess, sys, pathlib

REQS = ["numpy", "scipy", "numba", "matplotlib", "spiceypy"]
for pkg in REQS:
    try:
        __import__(pkg)
    except ImportError:
        subprocess.check_call(
            [sys.executable, "-m", "pip", "install", "--quiet", pkg])

ROOT = pathlib.Path.cwd().resolve().parents[1]   # repo root
sys.path.insert(0, str(ROOT))
from lunar.data_fetch import ensure_pds, ensure_spice
ensure_pds("apollo_hfe")     # fetch Apollo HFE tables from NASA/PDS
ensure_spice("de440")        # fetch DE440 kernel from NAIF
\end{lstlisting}

A reviewer therefore reproduces every headline figure in this book in
a single click, without manual data provisioning.  The full four-site
validation executes end-to-end on a 2024-class laptop in under
\SI{15}{\minute}.

\begin{milestone}
\textbf{Reproducibility acceptance:} a clean \code{git clone} of
\repo, followed by \emph{Run All} on
\code{01\_apollo\_validation.ipynb}, regenerates
\cref{fig:a15-timeseries,fig:a15-meanT,fig:a17-meanT,fig:change4-val,fig:chaste-val}
with RMSE values matching \cref{tab:phase1-rmse} to
\SI{1e-3}{\kelvin}.
\end{milestone}

\section{Summary of claims}

In one sentence: \emph{this thesis defends the claim that a single
1-D Crank-Nicolson solver, equipped with a pixel-level topographic
boundary condition and an ice-coupled conductivity closure, can
reproduce every in-situ lunar subsurface thermal record currently
published, and can be extended to a polar map without retuning the
physics kernel.}  The remaining chapters of this book lay out the
evidence for each half of that claim in sequence.
